\documentclass[12pt,a4paper]{article}

% Packages for math, code, and formatting
\usepackage{amsmath, amssymb, amsthm}
\usepackage{algorithm}
\usepackage{algpseudocode}
\usepackage{listings}
\usepackage{xcolor}

% Code style for C++
\lstdefinestyle{cpp}{
  language=C++,
  basicstyle=\ttfamily\small,
  keywordstyle=\color{blue},
  commentstyle=\color{green!50!black},
  stringstyle=\color{orange},
  numbers=left,
  numberstyle=\tiny,
  stepnumber=1,
  numbersep=5pt,
  frame=single,
  breaklines=true,
  tabsize=2,
  showstringspaces=false
}

% Theorem environments (for invariants, correctness proofs, etc.)
\newtheorem{theorem}{Teorema}
\newtheorem{lemma}{Lema}
\newtheorem{invariant}{Invariante}
\newtheorem{proposition}{Proposição}

\title{Relatório MC521 -  Problemas de Fluxo}
\author{Sofia Valverde Villas Bôas - 252974}

\begin{document}
\maketitle

\section*{A - Download Speed}

Nesse problema, temos uma rede consistindo de $n$ computadores e $m$ conexões. 
Cada conexão tem uma velocidade de download associada. Sendo $1$ o servidor e $n$ o de 
Kotivalo, queremos saber a velocidade máxima de download que Kotivalo pode obter para baixar um arquivo do
servidor até o seu computador.

\subsection*{Solução}

A ideia aqui foi modelar o problema como um grafo e, com isso, obter o fluxo máximo deste grafo.
Isso porque, podemos interpretar a velocidade de download de cada conexão como a capacidade de uma aresta (o que está limitando
a velocidade que Kotivalo pode obter). Assim, o fluxo máximo do grafo é a velocidade máxima de download que é pedida no output.

Portanto, bastou usar o algoritmo de Edmonds-Karp, inspirado no do material da aula, para computar isto.
Resumidamente, o algoritmo se baseia em obter caminhos aumentantes através de uma BFS,
e ir aumentando o fluxo do grafo até que não seja mais possível obter um caminho no grafo residual. 

Por fim, bastou imprimir o fluxo máximo, que é a resposta do problema.
\section*{C - School Dance}

Temos $n$ meninos e $m$ meninas. Um par de dança é formado por um menino e uma menina.
Há $k$ pares possíveis dados no enunciado e o problema pede para obtermos o número máximo de pares de dança que podem ser formados,
e quais são esses pares.

\subsection*{Solução}

Neste problema, da mesma forma que no anterior, podemos o modelar como um grafo e, dada a definição do enunciado,
sabemos que ele é bipartido. O que queremos obter então nada mais é do que o emparelhamento máximo do grafo, já que
não podemos ter mais de um par de dança para um mesmo menino ou para uma mesma menina. Ou seja, não podemos ter um 
mesmo vértice do grafo aparecendo em mais de uma aresta do emparelhamento.

A solução, então, consistiu em criar uma fonte que se liga a cada menino, e um sorvedouro que se liga a cada menina.
Para garantir que cada menino ou menina só possa ser emparelhado com um outro vértice, basta colocar capacidade 1 em cada
aresta que liga a fonte a um menino e em cada aresta que liga uma menina ao sorvedouro. Isso porque, mesmo que um menino ou
menina tenha mais de um par possível, o algoritmo de fluxo máximo só vai conseguir enviar 1 unidade de fluxo através do
vértice correspondente a esse menino ou menina, já que da aresta que liga a fonte ao menino, ou a menina ao sorvedouro,
só pode passar fluxo 1 no máximo. Assim, apenas uma das arestas que saem (ou chegam) desse vértice vai conseguir ter
fluxo igual a 1, e as demais ficam com fluxo 0. As arestas que ligam meninos a meninas também recebem capacidade 1.

Por fim, para obter os pares de dança, foi necessário iterar para cada menino e verificar
se alguma das arestas que saem do vértice correspondente a ele tem fluxo 1, o que em código
se traduziu em verificar se $cap[v][i] > 0$ para algum $v$ vizinho de $i$ (sendo $i$ o vértice do menino). Isto é,
se no grafo residual, o fluxo que podemos desfazer é de $1$. Isso funciona porque, no grafo residual, a aresta reversa
(de $v$ para $i$) começa com capacidade 0 e é incrementada exatamente pela quantidade de fluxo que foi enviada de
$i$ para $v$. Logo, $cap[v][i] > 0$ indica que 1 unidade de fluxo foi enviada por essa aresta, ou seja, que o par
$(i, v)$ faz parte do emparelhamento. Neste caso, imprimimos o par correspondente e passamos para o próximo menino.
\end{document}